\documentclass[11pt]{article}

\usepackage[T1]{fontenc}
\usepackage[utf8]{inputenc}
\usepackage{lmodern}
\usepackage{geometry}
\usepackage{setspace}
\usepackage{amsmath,amssymb,amsthm}
\usepackage{hyperref}
\usepackage{enumitem}
\usepackage{authblk}
\usepackage{orcidlink}
\usepackage{graphicx}
\usepackage{booktabs}
\usepackage{longtable}
\usepackage{array}
\usepackage{caption}
\usepackage[section]{placeins}
\usepackage{float}
\usepackage{adjustbox}
\usepackage{microtype}
\usepackage{xcolor}

\setlength{\textfloatsep}{10pt plus 2pt minus 2pt}
\setlength{\floatsep}{10pt plus 2pt minus 2pt}
\setlength{\intextsep}{10pt plus 2pt minus 2pt}
\renewcommand{\topfraction}{0.9}
\renewcommand{\bottomfraction}{0.8}
\renewcommand{\textfraction}{0.07}
\renewcommand{\floatpagefraction}{0.8}

\geometry{margin=1in}
\setstretch{1.15}

\hypersetup{
    colorlinks=true,
    linkcolor=blue,
    urlcolor=blue,
    citecolor=blue
}

\newcommand{\priorpreprintdoi}{10.5281/zenodo.19827222}
\newcommand{\terminalrows}{935{,}111}
\newcommand{\terminalmin}{28.499547\mathrm{B}}
\newcommand{\terminalmax}{30.610043\mathrm{B}}
\newcommand{\primarypackets}{4{,}676}
\newcommand{\projectionp}{0.1417165669}
\newcommand{\piminusthree}{0.1415926536}

\newcommand{\optionalfigure}[3]{%
\begin{figure}[H]
    \centering
    \IfFileExists{#1}{%
        \includegraphics[width=0.98\textwidth,
            height=0.80\textheight,keepaspectratio]{#1}%
    }{%
        \fbox{\parbox[c][2.25in][c]{0.88\textwidth}{%
            \centering Figure placeholder: \texttt{\detokenize{#1}}}}%
    }
    \caption{#2}
    \label{#3}
\end{figure}}

\title{\textbf{Terminal Altitude Architecture and Post-Selection Artifacts}\\
\textbf{in Adjacent Zeta-Zero Spacing}}

\author[1]{Jeffery Huckstead\,\orcidlink{0009-0007-0234-2177}}
\affil[1]{Independent Researcher}

\date{\today\\
Replacement manuscript draft for the previously released phase-locking preprint\\
Original preprint DOI: \href{https://doi.org/\priorpreprintdoi}{\priorpreprintdoi}}

\begin{document}

\maketitle

\begin{abstract}
We re-examine extreme adjacent-spacing compression packets among Riemann zeta zeros in the terminal LMFDB band using a seam-stitched board of \terminalrows{} rows spanning approximately \(\terminalmin\) to \(\terminalmax\). Earlier exploratory projections suggested two positive structures: a narrow first-logarithmic-harmonic \(h_1\) corridor and a discrete row-sector structure centered on modulus \(1009\). Subsequent search-adjusted validation audits show that neither interpretation survives as an explanatory mechanism. The circular resultant \(R\) is reproduced or exceeded by packet-count-matched search-adjusted synthetic and permutation controls; the apparent \(h_1\)-specific Arc80 corridor is absorbed by rank-altitude and altitude-matched controls; and the \(M=1009\) sector signal fails a dense \(M=800,\ldots,1200\) search-adjusted permutation null.

The surviving empirical object is a multi-component terminal altitude architecture in the yield rate of highly compressed adjacent-spacing packets. A primary reconstructed packet population contains \(\primarypackets\) extreme compression packets. Its altitude distribution is strongly non-uniform after normalization by board-row density, with descriptive Gaussian-mixture summaries favoring a multi-component altitude model over a single-component model. Altitude-only surrogate packet sets reproduce the previously observed \(h_1\) Arc80 corridor, showing that \(h_1\) is a descriptive projection of altitude localization rather than a coordinate-specific phase resonance. The resulting claim is finite-horizon and deliberately narrower: extreme compression packets form a terminal altitude architecture, while the earlier \(h_1\) and modulus-\(1009\) structures are examples of visually compelling post-selection artifacts in high-dimensional exploratory analysis of zeta-zero spacings.
\end{abstract}

\noindent\textbf{Keywords:} Riemann zeta zeros; adjacent zero spacing; extreme compression; terminal altitude architecture; post-selection inference; look-elsewhere effect; seam stitching; random matrix statistics; empirical number theory

\section{Introduction}

The local statistics of the nontrivial zeros of the Riemann zeta function have long been compared with random matrix ensembles, especially the Gaussian Unitary Ensemble (GUE) after unfolding \cite{montgomery1973,odlyzko1987,mehta2004}. These comparisons primarily concern bulk spacing statistics and correlation laws. The present work concerns a narrower empirical object: rare extreme compression events in adjacent zero spacing on a finite observed horizon.

An earlier version of this study reported that extreme compression packets appeared to organize in the first logarithmic harmonic
\[
h_1(T)=\left\{\frac{\log(T/2\pi)}{\pi}\right\}.
\]
A continuation study then reported a possible row-sector structure centered on modulus \(1009\). Both interpretations were plausible exploratory descriptions, and both were supported by local nulls conditional on the chosen coordinate or modulus family. A blind-review critique correctly identified the central remaining weakness: the nulls did not fully simulate the act of searching over coordinate families, harmonic families, moduli, window sizes, and tail thresholds.

This replacement manuscript reports the result of the subsequent search-adjusted validation campaign. The original \(h_1\)-specific interpretation does not survive. The \(M=1009\) row-sector interpretation does not survive. The underlying extreme-compression packet field, however, does survive. The data now support a different claim: the packets form a multi-component altitude architecture in the terminal band, and the old \(h_1\) corridor is explained as a projection of that altitude architecture through a slowly varying logarithmic coordinate.

The paper therefore has two goals. First, it reports the surviving empirical structure: a terminal altitude architecture in the yield rate of highly compressed adjacent-spacing packets. Second, it documents the failure modes of two attractive but ultimately non-surviving projections: the continuous \(h_1\) phase projection and the discrete modulus-\(1009\) sector projection. The result is both an empirical finding and a methodological caution about post-selection inference in computational number theory.

\subsection{What this manuscript replaces}

The original preprint framed the terminal packet field as evidence for a coordinate-specific \(h_1\) phase-locking phenomenon. That language is no longer maintained. In the present manuscript:
\begin{itemize}
    \item the \(h_1\) coordinate is retained only as a descriptive projection;
    \item the circular resultant \(R\) is demoted to a diagnostic statistic;
    \item the narrow \(h_1\) Arc80 corridor is explained by altitude localization;
    \item the modulus-\(1009\) row-sector result is rejected under dense search-adjusted controls;
    \item the surviving positive claim is terminal altitude architecture of extreme compression packets.
\end{itemize}

\subsection{Scope}

All claims are finite-horizon empirical claims on the loaded terminal board. No asymptotic theorem is claimed. No proof or disproof of the Riemann Hypothesis is claimed. The work assumes the tabulated zeros used in the computation are zeros on the critical line and studies their adjacent spacings as an empirical sequence.

\section{Data and seam-stitched board}

\subsection{Source horizon}

The terminal board used in the validation campaign contains \(\terminalrows\) rows and spans approximately
\[
T = \terminalmin \quad \text{to} \quad T = \terminalmax.
\]
The data are drawn from the available LMFDB zero tables, supplemented conceptually by the classical Odlyzko zero-table framework \cite{lmfdb,odlyzko_tables}. The working filenames used during the computational campaign sometimes contain the string \texttt{31p5B}; this was a downloader ceiling, not a claim that source data extend to \(31.5\mathrm{B}\). The actual terminal ordinate in the loaded source sequence is approximately \(30.610\mathrm{B}\).

\subsection{Seam stitching}

The board is seam stitched: consecutive zero ordinates are decoded in shard order and cross-shard adjacent gaps are retained. This matters because discarding cross-shard gaps would create artificial seams in any local spacing statistic. The seam-preserving design is therefore part of the empirical instrument, not a cosmetic preprocessing step.

\subsection{Terminal row metric}

The terminal packet audits use row-level spacing ratios derived from fitted and pure spacing spines. The primary metric in the final validation chain is \texttt{Gap\_to\_Fitted\_Ratio}; \texttt{Gap\_to\_Pure\_Ratio} is retained as a diagnostic comparator. Extreme packets are extracted by thresholding the row-level score distribution at high quantiles and grouping adjacent selected rows into connected packets.

The primary row-level reconstruction used in the validation chain selects the upper tail of \texttt{Gap\_to\_Fitted\_Ratio} at
\[
\tau=0.995,
\]
yielding \(\primarypackets\) packets. This is close to the original terminal packet count used in the earlier phase-locking analysis, and it avoids the one-packet degeneracy encountered by an initial rolling-window proxy benchmark.

\section{Methods}

\subsection{Packet extraction}

For a row-level score sequence \(X_k\), threshold \(\tau\), and tail direction \(d\), define the selected row set
\[
\mathcal{E}_{\tau,d}
=
\begin{cases}
\{k:X_k\ge q_\tau(X)\}, & d=\mathrm{ge},\\[4pt]
\{k:X_k\le q_{1-\tau}(X)\}, & d=\mathrm{le}.
\end{cases}
\]
A packet is a connected component of \(\mathcal{E}_{\tau,d}\) in row-index order. The packet center is chosen as the most extreme row within the component.

The primary population in this manuscript uses \(\texttt{Gap\_to\_Fitted\_Ratio}\), \(d=\mathrm{ge}\), and \(\tau=0.995\). Sensitivity grids over nearby \(\tau\) values and over the pure-ratio score are used to distinguish stable architecture from metric-specific endpoint behavior.

\subsection{Circular statistics}

For packet altitudes \(T_1,\ldots,T_N\), the earlier preprint studied the phases
\[
h_m(T_j)=\left\{m\frac{\log(T_j/2\pi)}{\pi}\right\}.
\]
The circular resultant is
\[
R=\left|\frac{1}{N}\sum_{j=1}^N e^{2\pi i h_m(T_j)}\right|.
\]
For \(q=0.80\), Arc80 is the shortest circular arc covering \(80\%\) of the packet phases.

The validation campaign shows that \(R\) is not a robust positive statistic under search-adjusted nulls. Arc80 remains useful, but the \(h_1\) Arc80 corridor is explained by altitude localization rather than by \(h_1\)-specific resonance.

\subsection{Packet yield rate}

A central correction in this replacement manuscript is the use of packet yield rate rather than raw packet count. For an altitude bin \(B\), define
\[
Y(B)=\frac{\#\{\text{packet centers in }B\}}{\#\{\text{board rows in }B\}}.
\]
This avoids the denominator trap: zero density varies with altitude, so raw packet counts alone cannot distinguish true packet overrepresentation from a larger row population.

\subsection{Altitude-only projection null}

To test whether the \(h_1\) Arc80 corridor is explained by altitude localization, we generate surrogate packet sets that preserve the packet altitude histogram. The terminal altitude range is divided into bins; within each bin, the surrogate draws the same number of rows as the real packet field contributes from that bin. The surrogate packet altitudes are then mapped through \(h_1\), and Arc80 is recomputed.

If these altitude-only surrogates reproduce the real \(h_1\) Arc80, then the \(h_1\) corridor is explained as a projection of altitude localization.

\subsection{Search-adjusted harmonic null}

The \(h_1\) claim is vulnerable to a look-elsewhere effect over harmonics. The validation campaign therefore allows synthetic and permutation nulls to scan
\[
h_1,h_2,h_3,h_4,h_5
\]
and retain the best circular statistic. This compares the real result to the null's best available harmonic, not merely to the null's \(h_1\) value.

\subsection{Coordinate-specificity audit}

To test whether the Arc80 corridor belongs specifically to \(h_1\), the validation campaign compares \(h_1\) against:
\begin{itemize}
    \item nearby frequency perturbations \(h_{1,\epsilon}(T)=\{(1+\epsilon)\log(T/2\pi)/\pi\}\);
    \item higher harmonics \(h_m\);
    \item monotone altitude controls, including linear \(T\), \(\sqrt{T}\), \(\log T\), and rank-\(T\);
    \item random phase controls.
\end{itemize}
The key normalized statistic is
\[
\mathrm{Arc80}_{\mathrm{norm}}=
\frac{\mathrm{Arc80}(\text{packet phases})}
     {\mathrm{Arc80}(\text{all board rows})}.
\]
This prevents coordinates with small total phase motion over the terminal band from appearing artificially strong.

\subsection{Dense modulus search-adjusted null}

For the row-sector claim, the validation campaign scans all moduli
\[
M=800,801,\ldots,1200
\]
under a base-6 sectorization and computes \(S_{\mathrm{sector}}(M)\). Each permutation null shuffles tail signs, scans the same full modulus range, and records
\[
\max_{800\le M\le 1200} S_{\mathrm{sector}}^{\mathrm{null}}(M).
\]
This is the correct look-elsewhere comparison for the modulus claim.

\section{Validation chronology}

Table~\ref{tab:validation-chronology} summarizes the validation sequence.

\begin{table}[H]
\centering
\small
\begin{adjustbox}{max width=\textwidth}
\begin{tabular}{p{0.12\textwidth}p{0.32\textwidth}p{0.31\textwidth}p{0.18\textwidth}}
\toprule
\textbf{Audit} & \textbf{Question} & \textbf{Result} & \textbf{Role} \\
\midrule
P83B & Does \(h_1\) concentration survive packet-count-matched search-adjusted synthetic controls? & Split result: \(R\) fails; Arc80 survives packet-count controls. & Demotes \(R\); isolates Arc80 for further testing. \\
P84 & Is the surviving Arc80 corridor \(h_1\)-specific? & No. Rank/altitude controls and block-matched nulls absorb \(h_1\) specificity. & Rejects coordinate-specific \(h_1\) interpretation. \\
P85 & What explains the \(h_1\) corridor if \(h_1\) is not the mechanism? & Altitude-only surrogates reproduce \(h_1\) Arc80; packet field has multi-component altitude architecture. & Establishes replacement mechanism. \\
P85R & Is there extra sub-block microstructure beyond macro altitude bands? & No strong support after repaired sub-block null. & Keeps replacement mechanism macro-altitude, not micro-fractal. \\
P86 & Does modulus \(1009\) survive dense search-adjusted modulus controls? & No. No modulus survives look-elsewhere correction. & Rejects row-sector engine claim. \\
\bottomrule
\end{tabular}
\end{adjustbox}
\caption{Validation chronology leading to the replacement interpretation.}
\label{tab:validation-chronology}
\end{table}

\section{Results}

\subsection{The packet field survives}

The first key result is that the packet object survives the interpretive collapse of the \(h_1\) coordinate. The P83B reconstruction extracts \(\primarypackets\) packets under the primary fitted-ratio tail rule. This establishes that the empirical object is not merely an artifact of a rolling-window proxy or a one-packet degeneracy.

\begin{table}[H]
\centering
\small
\begin{tabular}{lr}
\toprule
\textbf{Quantity} & \textbf{Value} \\
\midrule
Terminal rows & \terminalrows \\
Altitude range & \(\terminalmin\) to \(\terminalmax\) \\
Primary metric & \texttt{Gap\_to\_Fitted\_Ratio} \\
Primary direction & \texttt{ge} \\
Primary threshold & \(\tau=0.995\) \\
Primary packet count & \(\primarypackets\) \\
\bottomrule
\end{tabular}
\caption{Primary terminal packet reconstruction used in the validation chain.}
\label{tab:primary-packet-reconstruction}
\end{table}

\subsection{\texorpdfstring{\(R\)}{R} fails, Arc80 survives packet-count controls}

P83B compares the packet field against packet-count-matched synthetic and permutation controls, with the null allowed to scan harmonics \(h_1,\ldots,h_5\). The result is split. The circular resultant \(R\) does not separate from the search-adjusted null. In fact, the null routinely matches or exceeds the real resultant. The Arc80 corridor, however, is narrower than the packet-count-matched synthetic and permutation controls.

This split changes the role of the statistics. \(R\) is retained only as descriptive. Arc80 remains a real geometric feature, but it requires coordinate-specificity testing.

\subsection{The Arc80 corridor is not \texorpdfstring{\(h_1\)}{h1}-specific}

P84 tests whether the surviving Arc80 corridor belongs specifically to \(h_1\). It does not. The normalized \(h_1\) Arc80 is outperformed by rank-altitude controls, and block-matched altitude nulls absorb the apparent \(h_1\)-specific effect. The epsilon perturbation surface is effectively flat, indicating that \(h_1\) is not a sharp local optimum in the nearby log-phase family.

The conclusion is not that the packet field is random. It is that \(h_1\) is not the explanatory coordinate.

\subsection{Altitude-only surrogates reproduce the \texorpdfstring{\(h_1\)}{h1} corridor}

P85 provides the replacement explanation. The real \(h_1\) Arc80 for the primary packet field is approximately
\[
\mathrm{Arc80}_{h_1}=0.017454.
\]
Altitude-only surrogate packets, drawn to preserve the same altitude-bin distribution, have mean
\[
\mathrm{Arc80}_{h_1}^{\mathrm{surr}}\approx 0.017469.
\]
The empirical probability that the surrogate Arc80 is at least as tight as the real value is
\[
p\approx \projectionp.
\]
Thus the old \(h_1\) corridor is reproduced by altitude localization alone. The \(h_1\) coordinate is a descriptive projection of the packet altitude distribution.

The numerical proximity
\[
\projectionp \approx \pi-3
\]
is recorded as a future-hypothesis marker only. It is not used as evidence in the present statistical interpretation.

\begin{table}[H]
\centering
\small
\begin{tabular}{lr}
\toprule
\textbf{Projection statistic} & \textbf{Value} \\
\midrule
Real \(h_1\) Arc80 & \(0.017454\) \\
Altitude-only surrogate mean \(h_1\) Arc80 & \(0.017469\) \\
Empirical probability surrogate \(\le\) real & \(\projectionp\) \\
\(\pi-3\) reference & \(\piminusthree\) \\
Interpretation & \(h_1\) Arc80 explained by altitude projection \\
\bottomrule
\end{tabular}
\caption{Projection explanation of the old \(h_1\) Arc80 corridor. The surrogate preserves packet altitude distribution but randomizes row choice within altitude bins.}
\label{tab:h1-projection}
\end{table}

\subsection{Altitude architecture of the packet field}

The packet yield rate is not uniform across altitude. P85 computes
\[
Y(B)=\frac{\#\{\text{packets in altitude bin }B\}}
          {\#\{\text{board rows in altitude bin }B\}},
\]
avoiding raw-count bias from changing row density. The yield-rate maps show broad altitude bands of elevated and suppressed packet frequency.

A Gaussian-mixture descriptive summary favors a multi-component model. In the primary run, \(K=4\) is selected by BIC, with a BIC improvement over \(K=1\) of approximately
\[
2823.78.
\]
This is not interpreted as a theorem that exactly four physical bands exist. It is interpreted as strong descriptive evidence that a single unimodal altitude model is inadequate.

\begin{table}[H]
\centering
\small
\begin{tabular}{lr}
\toprule
\textbf{Altitude-architecture statistic} & \textbf{Value} \\
\midrule
Primary packet count & \(\primarypackets\) \\
Best descriptive GMM \(K\) & \(4\) \\
BIC improvement \(K=1\) minus best \(K\) & \(2823.78\) \\
Sub-block microstructure support after macro-block matching & Not supported \\
Interpretation & Macro-scale terminal altitude architecture \\
\bottomrule
\end{tabular}
\caption{Descriptive altitude-architecture summary.}
\label{tab:altitude-architecture}
\end{table}

\subsection{No strong evidence for extra sub-block micro-clustering}

The original P85 block-null statistic was degenerate because the null preserved the same block counts used by the test statistic. P85R repaired this by matching macro-block counts while testing finer micro-bin and nearest-neighbor statistics. The repaired test does not find strong evidence for extra fine-scale clustering beyond the macro altitude architecture.

Representative repaired sub-block \(p\)-values are:
\[
p_{\mathrm{Gini}}\approx 0.180,\quad
p_{\max}\approx 0.315,\quad
p_{\mathrm{NN,mean}}\approx 0.244,\quad
p_{\mathrm{NN,median}}\approx 0.353.
\]
This simplifies the replacement interpretation: the surviving object is macro-altitude architecture, not confirmed fractal or micro-scale clustering.

\subsection{Tail-polarity altitude asymmetry}

The primary fitted-ratio upper and lower tails are not identical in altitude distribution. In the primary \(\tau=0.995\) fitted-ratio comparison, the two tails have a Kolmogorov--Smirnov-type distance of approximately
\[
0.168
\]
and a Jensen--Shannon divergence of approximately
\[
0.0496.
\]
This supports a descriptive tail-polarity asymmetry in altitude. The effect is not stated as a causal thermodynamic mechanism. It is a finite-horizon distributional asymmetry between the two fitted-ratio tails.

\subsection{The modulus-\texorpdfstring{\(1009\)}{1009} row-sector claim fails}

P86 scans all moduli \(M=800,\ldots,1200\). Each permutation null is allowed to scan the same modulus range and record its maximum sector amplitude. Under this search-adjusted null, neither \(M=1009\) nor any real best modulus survives.

Across the four seam-scale windows \(W=900,1000,1100,1300\), the number of configurations in which \(M=1009\) survives the search-adjusted null is
\[
0.
\]
The number of configurations in which the best real modulus survives the search-adjusted null is also
\[
0.
\]
Thus the row-sector continuation paper's positive \(M=1009\) interpretation is rejected as a search-hardened discovery.

\begin{table}[H]
\centering
\small
\begin{tabular}{lrrrr}
\toprule
\textbf{Window} & \textbf{Best real modulus} & \textbf{Rank of 1009} & \textbf{\(p_{\mathrm{best}}\)} & \textbf{\(p_{1009}\)} \\
\midrule
900  & 1133 & 48  & 0.0798 & 1.0000 \\
1000 & 1087 & 151 & 0.2355 & 1.0000 \\
1100 & 1133 & 143 & 1.0000 & 1.0000 \\
1300 & 1022 & 58  & 0.9381 & 1.0000 \\
\bottomrule
\end{tabular}
\caption{Dense modulus validation summary. The search-adjusted null scans \(M=800,\ldots,1200\) on every permutation. Neither \(M=1009\) nor the best real modulus survives.}
\label{tab:modulus-failure}
\end{table}

\section{Claim transition ledger}

Table~\ref{tab:claim-transition} summarizes the shift from the earlier preprint interpretation to the present replacement interpretation.

\begin{table}[H]
\centering
\small
\begin{adjustbox}{max width=\textwidth}
\begin{tabular}{p{0.30\textwidth}p{0.33\textwidth}p{0.30\textwidth}}
\toprule
\textbf{Old claim} & \textbf{Validation result} & \textbf{Revised claim} \\
\midrule
Extreme compression packets phase-lock to \(h_1\). &
P83B and P84 reject \(R\) and \(h_1\)-specificity under search-adjusted and coordinate-control nulls. &
Rejected as a coordinate-specific phase-locking resonance. \\
\midrule
The \(h_1\) Arc80 corridor is intrinsically \(h_1\)-specific. &
P85/P85R show that altitude-only surrogates reproduce the \(h_1\) Arc80 corridor. &
\(h_1\) Arc80 is a descriptive projection of altitude-localized packet architecture. \\
\midrule
The packet object disappears if \(h_1\) fails. &
P85 shows packet yield-rate maps and multi-component altitude structure. &
The core object survives as terminal altitude architecture of highly compressed adjacent-spacing packets. \\
\midrule
\(M=1009\) is a row-sector engine. &
P86 dense modulus scan finds no search-adjusted survival. &
Rejected as a look-elsewhere artifact in sparse tail-sector scanning. \\
\bottomrule
\end{tabular}
\end{adjustbox}
\caption{Claim transition ledger from the earlier preprint interpretation to the present replacement interpretation.}
\label{tab:claim-transition}
\end{table}

\section{Discussion}

The validation campaign changes the interpretation of the original discovery but does not erase the empirical object. The earlier \(h_1\) analysis detected a real pattern because the packets were not uniformly distributed in altitude. The mistake was to assign the explanatory role to \(h_1\) itself. The revised interpretation is more conservative: \(h_1\) is a projection through which altitude localization became visible.

The modulus-\(1009\) claim fails more completely. Once the null is allowed to search the same dense modulus range, the apparent sector engine disappears into the look-elsewhere envelope. Unlike the \(h_1\) corridor, which is explained as a projection of a surviving altitude architecture, the modulus-\(1009\) interpretation does not currently have a surviving positive replacement. It remains valuable as a methodological warning.

The central surviving result is therefore the altitude architecture of extreme compression packets. This is a finite-horizon empirical structure. It may reflect nonstationarity in the terminal segment, long-range fluctuation in the observed zero sequence, or a more subtle interaction between local spacing tails and altitude. The present paper does not identify a theoretical mechanism for the altitude bands. It establishes that the bands exist in the loaded terminal board and that they explain the earlier \(h_1\) projection.

\subsection{Relation to GUE statistics}

The results do not claim a breakdown of GUE bulk statistics. The packet field lies in the rare-event tail of adjacent spacing behavior. The observation is that the extreme compression packets are not uniformly distributed over the terminal board. This is compatible with bulk GUE-like behavior while still allowing finite-horizon tail architecture.

\subsection{Lessons for empirical number theory}

This study demonstrates a general methodological point. Conditional nulls can be misleading after a large exploratory search. A null that tests a chosen coordinate after the coordinate has already been selected does not address the look-elsewhere effect. For empirical studies of zeta zeros, phase coordinates, modulus scans, and geometric projections should be evaluated using nulls that are allowed to search the same family of alternatives.

\section{Limitations}

The analysis is finite-horizon. The terminal board ends near \(30.610\mathrm{B}\), and no claim is made beyond the loaded source range. The altitude architecture may change, vanish, or sharpen at greater heights.

The packet extractor is empirical. Although it has been stress-tested across nearby thresholds and metrics, it is not derived from a theorem about zeta zeros. The GMM component count is descriptive, not a proof of exactly four physical bands.

The \(h_1\) projection explanation is strong on the observed board, but it does not imply that all future \(h_1\)-like behavior is an artifact. It means the \(h_1\) behavior measured in the present terminal packet field is explained by altitude localization.

The modulus-\(1009\) failure is specific to the tested row-sector pipeline and dense modulus search. It rejects the positive claim made by the continuation paper but does not rule out all possible arithmetic sectorizations under different, pre-specified theoretical frameworks.

\section{Conclusion}

The original \(h_1\) phase-locking interpretation and the \(M=1009\) row-sector interpretation do not survive search-adjusted validation. The surviving structure is a terminal altitude architecture in highly compressed adjacent-spacing packets. The packet field is non-uniform over altitude, has multi-component descriptive structure, and explains the previously observed \(h_1\) Arc80 corridor through altitude projection.

The revised conclusion is therefore narrower but stronger: the terminal band contains a finite-horizon architecture of extreme compression packets, and this architecture can cast misleading shadows in both continuous logarithmic phase coordinates and discrete modulus coordinates. Search-adjusted validation separates the surviving object from its projections.

\appendix

\section{Audit glossary}

\begin{longtable}{p{0.16\textwidth}p{0.78\textwidth}}
\toprule
\textbf{Audit} & \textbf{Role} \\
\midrule
P83B & Packet-count-matched search-adjusted harmonic benchmark. Demotes \(R\), preserves Arc80 as a statistic requiring coordinate testing. \\
P84 & Coordinate-specificity audit. Rejects \(h_1\)-specific explanation of the Arc80 corridor. \\
P85 & Altitude architecture audit. Establishes packet yield-rate altitude structure and shows altitude-only surrogates reproduce \(h_1\) Arc80. \\
P85R & Cleanup audit. Repairs the degenerate block-null and produces the claim transition ledger. \\
P86 & Dense modulus search-adjusted null. Rejects \(M=1009\) as a search-hardened sector-engine discovery. \\
P87 & Packaging scaffold. Consolidates the surviving and rejected claims into this replacement manuscript structure. \\
\bottomrule
\end{longtable}

\section{Notation}

\begin{table}[H]
\centering
\small
\begin{tabular}{ll}
\toprule
\textbf{Symbol} & \textbf{Meaning} \\
\midrule
\(\gamma_n\) & ordinate of the \(n\)-th nontrivial zero on the critical line \\
\(g_n\) & adjacent-zero gap \(\gamma_{n+1}-\gamma_n\) \\
\(T\) & representative altitude \\
\(h_m(T)\) & logarithmic harmonic phase coordinate \\
\(R\) & circular resultant \\
Arc80 & shortest circular arc covering \(80\%\) of packet phases \\
\(Y(B)\) & packet yield rate in altitude bin \(B\) \\
\(\tau\) & row-level tail threshold quantile \\
\bottomrule
\end{tabular}
\caption{Principal notation.}
\end{table}

\section*{Data and Artifact Availability}

The pre-validation preprint was deposited under DOI
\href{https://doi.org/\priorpreprintdoi}{\priorpreprintdoi}. The validation campaign generates the replacement claim ledger and manuscript scaffold in the ForgeV16c audit chain, including outputs from P83B, P84, P85, P85R, P86, and P87. The source zero ordinate data are attributed to LMFDB and Odlyzko zero tables \cite{lmfdb,odlyzko_tables}.

\section*{Use of AI Tools}

Generative AI tools were used for drafting assistance, code critique, adversarial claim review, and organization of the validation ledger. All research design decisions, computational execution choices, interpretation, and final manuscript responsibility remain with Jeffery Huckstead.

\section*{Author Contributions}

Jeffery Huckstead conceived the research program, assembled and audited the seam-stitched board, executed the computational validation campaign, interpreted the results, and wrote the manuscript.

\bibliographystyle{plain}

\begin{thebibliography}{9}

\bibitem{montgomery1973}
H.~L. Montgomery,
\emph{The pair correlation of zeros of the zeta function},
in \emph{Analytic Number Theory}, Proc.\ Sympos.\ Pure Math., vol.~24,
Amer.\ Math.\ Soc., Providence, RI, 1973, pp.~181--193.

\bibitem{odlyzko1987}
A.~M. Odlyzko,
\emph{On the distribution of spacings between zeros of the zeta function},
Math.\ Comp.\ \textbf{48} (1987), no.~177, 273--308.

\bibitem{mehta2004}
M.~L. Mehta,
\emph{Random Matrices},
3rd ed., Elsevier/Academic Press, Amsterdam, 2004.

\bibitem{rubinstein2005}
M.~Rubinstein,
\emph{Computational methods and experiments in analytic number theory},
in \emph{Recent Perspectives in Random Matrix Theory and Number Theory},
London Math.\ Soc.\ Lecture Note Ser., vol.~322,
Cambridge Univ.\ Press, Cambridge, 2005, pp.~425--506.

\bibitem{lmfdb}
The LMFDB Collaboration,
\emph{The \(L\)-functions and Modular Forms Database},
\url{https://www.lmfdb.org}, accessed 2024--2026.

\bibitem{odlyzko_tables}
A.~M. Odlyzko,
\emph{Tables of zeros of the Riemann zeta function},
available at \url{http://www.dtc.umn.edu/~odlyzko/zeta_tables/},
accessed 2024--2026.

\end{thebibliography}

\end{document}
